\section{Introduction}

Small RNAs complexed with Argonaute family proteins silence complementary
genomic regions either by degrading messenger RNAs or by directing
transcriptional gene silencing at chromatin \citep{holoch2015rna,wilson2013rnai}.
Transcriptional silencing generally requires that small RNAs recognize their
genomic targets via nascent long non-coding RNAs (lncRNAs). This leads to an
apparent paradox: chromatin regions destined for silencing must nevertheless be
transcribed so that they can be recognized by their cognate small RNAs.

Different organisms have evolved different solutions to this problem. In
fission yeast, siRNAs direct the deposition of histone H3 lysine~9
di- and tri-methylation (H3K9me2/3) on centromeric repeats, leading to
heterochromatin assembly by the HP1 protein Swi6 \citep{volpe2002rnai,hall2002heterochromatin}.
Mitotic phosphorylation of histone~H3 serine~10 evicts Swi6, providing a
transient window in which RNA polymerase~II can re-transcribe these repeats.
In Drosophila, the HP1 paralog Rhino binds H3K9me2/3 at piRNA clusters
\citep{lethomas2014rhino} and actively recruits dedicated variants of basal
transcription factors so that lncRNAs can be produced from heterochromatic
loci while mRNA transcription of the same transposons is suppressed
\citep{andersen2017heterochromatin}. In plants, the SHH1 reader recruits the
plant-specific RNA polymerase~IV to H3K9-methylated chromatin to license
lncRNA transcription for RNA-directed DNA methylation \citep{law2013shh1}.

In most of these systems the source and target of the small RNA coincide,
making it difficult to tease apart lncRNA transcription for small RNA
biogenesis from lncRNA transcription for effector recruitment. Programmed DNA
elimination in ciliates such as \emph{Tetrahymena thermophila} and
\emph{Paramecium tetraurelia} provides a particularly clean system because
small RNA source and target loci reside in physically separate nuclei
\citep{chalker2011ciliate}. Each ciliate cell contains a diploid germline
micronucleus (MIC) and a polyploid somatic macronucleus (MAC); during
conjugation the MIC undergoes meiosis and fertilization to form new MIC and
MAC copies, while the parental MAC is destroyed. DNA elimination takes place
in the developing new MAC, where tens of thousands of internal elimination
sequences (IESs) are excised to convert the approximately $200$\,Mb
micronuclear genome into the approximately $103$\,Mb macronuclear genome
\citep{cheng2010transposase,dubois2017piggymac}.

Three classes of lncRNAs coordinate this process, and they appear in three
different nuclei at distinct times. During meiotic prophase, bidirectional
transcription by RNA polymerase~II produces genome-wide MIC-lncRNAs in the
germline micronucleus \citep{mochizuki2004rpb3,chalker2001nongenic}; these
are diced into $\sim$29\,nt scan RNAs (scnRNAs) and loaded onto the Piwi
protein Twi1 \citep{mochizuki2002twi1,noto2010giw1}. The Twi1--scnRNA complex
then moves into the parental MAC, where the RNA helicase Ema1 promotes the
interaction between scnRNAs and nascent parental MAC lncRNAs (pMAC-lncRNAs)
\citep{aronica2008ema1}. This interaction triggers target-directed small RNA
degradation (TDSD), which eliminates MAC-matching scnRNAs and retains those
complementary to germline-limited sequences \citep{schoeberl2012biased}.
When the new MAC develops, the retained scnRNAs interact with nMAC-lncRNAs to
recruit Polycomb Repressive Complex~2 (PRC2), which catalyses both
H3K9me2/3 and H3K27me3 \citep{liu2004h3k9,liu2007rnai,frapporti2019ezl1} for
IES-specific heterochromatin assembly, followed by excision by domesticated
PiggyBac transposases \citep{cheng2010transposase,dubois2017piggymac}.

Despite the central role of pMAC-lncRNAs in TDSD, the regulators that enable
their transcription have remained unknown. Here we identify Ema2
(TTHERM\_00113330), a conjugation-specific SUMO E3 ligase with an atypical
SP-RING domain, as the first factor required for genome-wide pMAC-lncRNA
accumulation. Using knockout analysis, in~vitro binding assays, mass
spectrometry and targeted lysine-to-arginine mutants, we show that Ema2
interacts with the SUMO E2 enzyme Ubc9, catalyses SUMOylation of the
transcription elongation factor Spt6, and promotes the association of Spt6
and RNA polymerase~II with chromatin. These activities are necessary for
TDSD, for H3K27me3 deposition in the parental MAC, and ultimately for the
production of viable sexual progeny.

\paragraph{Contributions.}
\begin{itemize}
  \item We identify Ema2 as a conjugation-specific SUMO E3 ligase required
        for the accumulation of pMAC-lncRNAs, TDSD and programmed DNA
        elimination in \emph{T.~thermophila}.
  \item We show that Ema2 interacts directly with the SUMO E2 Ubc9 in vitro
        and is responsible for the majority of conjugation-stage SUMOylation,
        including SUMOylation of the transcription elongation factor Spt6.
  \item We demonstrate that Ema2 facilitates chromatin association of Spt6
        and RNA polymerase~II in the parental MAC at the mid-conjugation
        stage, offering a mechanism for how SUMOylation supports lncRNA
        transcription.
  \item By constructing and analysing a SUMOylation-defective Spt6 mutant, we
        show that Spt6 SUMOylation per se is dispensable for pMAC-lncRNA
        transcription but still contributes to efficient DNA elimination,
        indicating that Ema2 operates through additional SUMO targets.
\end{itemize}

\section{Related Work}

\paragraph{Small RNA-directed heterochromatin and lncRNA transcription.}
Several lineages use distinct transcriptional strategies to produce lncRNAs
that anchor small RNA-directed silencing. Centromeric transcription in fission
yeast is enabled passively by cell-cycle-dependent eviction of Swi6
\citep{volpe2002rnai,hall2002heterochromatin}, whereas Drosophila piRNA
clusters employ dedicated Rhino-containing transcription machinery that
exploits H3K9me2/3 marks \citep{lethomas2014rhino,andersen2017heterochromatin}.
In plants, SHH1 recruits Pol~IV at H3K9-methylated loci to drive siRNA
precursor transcription \citep{law2013shh1}. Our work adds a mechanistically
distinct solution: Ema2-directed SUMOylation acts \emph{upstream} of lncRNA
transcription in Tetrahymena, in a heterochromatin-independent manner.

\paragraph{SUMOylation in small RNA pathways.}
SUMO modifications are tightly linked to RNA-directed silencing across
eukaryotes. In fission yeast, SUMOylation of Swi6, Chp2 and Clr4 is
important for heterochromatin maintenance \citep{shin2005sumo}. In
\emph{C.~elegans}, SUMOylation of HDAC1 promotes Argonaute-directed
transcriptional silencing \citep{kim2021hdac1}. In Drosophila, the SUMO E3
ligase Su(var)2-10 couples piRNA target recognition to SetDB1 recruitment
and H3K9me3 deposition \citep{ninova2019suvar}. These examples all act
downstream of existing lncRNAs, whereas our data demonstrate that Ema2
enables lncRNA production itself.

\paragraph{SP-RING domain SUMO E3 ligases.}
Structural studies of Siz/PIAS family ligases in \emph{S.~cerevisiae}
(Siz1) and the Smc5/6-associated Mms21 have defined the zinc-coordinating
cysteine and histidine residues of the canonical SP-RING, and the accessory
residues that stabilise the RING--E2 interface
\citep{yunus2009siz1,duan2009mms21,hochstrasser2001spring}. The recognition
motif $\psi$KxE and its direct engagement by Ubc9 provide a baseline for
substrate selectivity \citep{sampson2001sumo}. Ema2 deviates from this
template in that the first zinc-binding cysteine is replaced by histidine
and several stabiliser positions appear to be missing, so we analyse its
activity biochemically and genetically rather than inferring it from sequence
alone.

\paragraph{SUMOylation and DNA elimination in ciliates.}
In \emph{Paramecium tetraurelia}, RNAi knockdown of \emph{UBA2} or \emph{SUMO}
was previously shown to block internal eliminated sequence excision, implying
a conserved role for the SUMO pathway in ciliate genome remodelling
\citep{matsuda2006sumo}. In \emph{Tetrahymena thermophila}, SUMOylation
increases dramatically during conjugation and is required for cell pairing
\citep{nasir2015sumo}. Our identification of Ema2 as a major conjugation-stage
SUMO E3 places a specific molecular activity within this pathway and links
it to pMAC-lncRNA transcription, TDSD and heterochromatin deposition.

\paragraph{TDSD and target-directed microRNA degradation.}
TDSD in Tetrahymena is conceptually reminiscent of target-directed microRNA
degradation and of the Ema1-dependent base pairing between scnRNAs and
pMAC-lncRNAs that selects IES-matching small RNAs for survival
\citep{aronica2008ema1,schoeberl2012biased}. By inactivating pMAC-lncRNA
transcription upstream of base pairing, Ema2 provides a genetically separable
entry point into this process.
